% !TEX root = ../main.tex
% !TEX program = xelatex
\title{融合多源预测的微网经济调度与线性优化研究}
\date{}
\maketitle
\begingroup
\setstretch{1.22}
\begin{abstract}
\noindent\hspace*{2em}随着我国碳达峰、碳中和目标的推进和分布式新能源的发展，光伏与储能在社区供能中的应用不断拓展。然而，光伏出力与用电需求存在时间错位，提前订电又面临供需和价格变化：订电不足可能产生高价紧急购电，订电过多则可能付费后无法充分利用。为降低微网购电费用，本文以10分钟为基本时段，建立统一的电量流与储能状态关系，研究不同信息条件下的购电及充放电策略。

\textbf{针对问题一}，采用电量流线性规划，在日初与日末储电量相等的条件下统筹全天购电与充放电，将有限储能留给更高价的时段。代表日最低购电费用为\textbf{35126.85元}，较峰谷策略降低8.23\%，表明全天统筹能够改善仅按价格门槛安排储能的效果。

\textbf{针对问题二}，将\textbf{近期加权准备与跨日线性规划}相结合。负载与光伏分别利用星期规律和近期历史预测，再根据近期缺口误差设置准备，平衡提前订电与五倍电价补购的损失。2025年2月至12月总费用为1358.38万元。与等权误差准备相比，普通电费增加约1.97万元，紧急电费减少约46.63万元，净节省\textbf{44.66万元}，说明适当增加普通购电准备能够减少更昂贵的临时补购。

\textbf{针对问题三}，采用\textbf{融合预测与日内校正滚动规划}，融合官方光伏预报与近期历史，依据当天已发生的负载偏差修正后续判断，并权衡调整费与补购费，更新尚未执行的订单。同一评价期总费用为1301.42万元。进一步设置相同模型下的信息对照：相较零点固定订单，日内反馈与改单节省\textbf{50.63万元}；在已有反馈与改单能力的基础上，后发官方光伏预报额外节省\textbf{6.95万元}，并使紧急购电量下降30.10\%。后两组期末库存相同，表明新增信息的收益并非来自多消耗剩余电量。

\textbf{针对问题四}，在未来真实电价未知的条件下，零点分支采用提升树预测与线性规划，日内分支采用岭回归预测与滚动线性规划；规划使用预测电价，结算采用实际电价。两套完整策略同期费用分别为1424.86万元和1369.59万元，后者降低3.88\%，反映供需判断、价格预测与订单调整共同作用下的节费效果。

逐时复核表明，方案满足供需平衡及储能容量、功率约束，费用可由实际订单和补购量重新计算。准备量检验与期末库存比较进一步说明了节费的条件与代价。本文将预测偏差、储能调节和信息更新与实际购电损失联系起来，形成了收益来源清晰、运行结果可核验的微网调控方案。
\end{abstract}
\noindent{\heiti\zihao{-4} 关键词：}{\songti\zihao{-4}\textbf{微网调度；储能优化；线性规划；滚动优化；电价预测}}
\par
\endgroup
